\chapter{PENDAHULUAN}
\section{Latar Belakang}


Tuliskan di sini \textbf{Pedoman Penulisan Latar Belakang:} Uraikan alasan pemilihan lokasi dan judul PKL jaringan komputer. Deskripsikan secara umum profil industri/perusahaan, urgensi ketersediaan infrastruktur jaringan komputer, dan permasalahan jaringan yang melatarbelakangi pengerjaan laporan ini.


\section{Tujuan dan Manfaat}


Tuliskan di sini \textbf{Pedoman Penulisan Tujuan dan Manfaat:} Jelaskan tujuan utama menganalisis jaringan komputer berjalan tersebut (misal: mengidentifikasi kelemahan sekuritas, menguji bandwidth, atau merancang topologi baru). Sebutkan manfaat praktisnya bagi perusahaan tempat PKL dan bagi pembaca.


\section{Metode Penelitian}


Tuliskan di sini \textbf{Pedoman Penulisan Metode Penelitian:} Deskripsikan metode yang Anda gunakan dalam mengumpulkan data jaringan, seperti observasi fisik perangkat keras jaringan, wawancara dengan network administrator setempat, pengujian Ping/Traceroute, atau studi literatur manual perangkat keras.


\section{Ruang Lingkup}


Tuliskan di sini \textbf{Pedoman Penulisan Ruang Lingkup:} Batasi ruang lingkup analisis jaringan Anda agar tidak terlalu luas (misalnya membatasi pada analisis jaringan LAN di lantai satu, konfigurasi router Mikrotik utama, sistem keamanan nirkabel/Wi-Fi, atau bandwidth management divisi tertentu saja).

